\section{Discussion}
\label{sec:discussion}

Three findings stand out. First, the distinction between counting settlements and counting people is not a technicality: a \PopulationMatchedPct\% settlement-match rate looks like a serious data gap, but the same frame represents \FrameCoveragePct\% of the census population, because the unmatched settlements are overwhelmingly small, informal, or education-system-only reference points rather than populous census-recognized centros poblados. Any accessibility analysis that reports coverage purely in settlement-count terms risks materially misstating how much of the population it actually describes.

Second, calibrating survey weights against a known auxiliary population total changed the substantive result, not just its precision: the calibrated mean access time (\MeanAccessTotal{} min) is \MeanAccessDiffHtCal{} minutes lower than the uncalibrated Horvitz--Thompson estimate, and department-level uncalibrated estimates were biased by roughly $-23\%$ to $+27\%$ relative to the calibrated ones. This is exactly the situation calibration is designed for: a probability sample with known, more-precise auxiliary totals at a finer geographic level than the original stratification. We report both figures throughout so that the correction remains auditable rather than a hidden adjustment.

Third, distinguishing ``poor but measured access'' from ``access we cannot yet measure'' matters as much as the headline mean. \CovNoTimeTotal\% of the frame population has no estimable route at all -- a data-availability statement -- while a separate \CovGtOneTwentyTotal\% has a measured time exceeding two hours. Collapsing these two categories, or silently excluding the first, would overstate how well the current data describe the hardest-to-reach places, which are disproportionately likely to be exactly where the no-time-estimate share is concentrated (Amazonas and rural Cusco).

The straight-line-vs-network comparison (median ratio \SvnMedian) reinforces why the whole exercise of network routing was necessary in the first place: a planner working from map distance alone would, for a typical settlement, underestimate true travel distance by roughly 70\%, and by far more for some settlements, with no constant correction available.
